\subsection*{一步一步建立直观}

\paragraph{1. 先为复数选择一个角度}

单位虚数 $i$ 位于单位圆的上端，主辐角为 $\pi/2$，但加上任意整数圈仍表示同一个点。右图位置只能显示点本身，不能区分绕了多少圈。复对数必须记录所选辐角，因此同一个非零复数会对应不同的对数值。



\paragraph{2. 把指数定义中的分支写出来}

若取主值，$\log i=i\pi/2$，于是 $i^i$ 按 $e^{i\log i}$ 计算，指数变成实数 $-\pi/2$。右图的半圈是几何参照，不是 $i^i$ 的输出位置；本窗口演示辐角与分支，实际幂值需要用指数定义计算。选定分支后，再给出一个确定数值才具有明确含义。



\paragraph{3. 每多记一圈，结果都会改变}

取其他对数分支时，$\log i=i(\pi/2+2k\pi)$，相应幂值为 $e^{-\pi/2-2k\pi}$。它们都是正实数，但彼此不同。因此“它是实数”并不意味着“没有分支选择”；给出一个数值答案前，必须明确采用主值还是多值约定。


